\documentclass{article}
\usepackage{amsthm}

\usepackage[utf8]{inputenc}     % pdfLaTeX
\usepackage[T1]{fontenc}
\usepackage[magyar]{babel}

\title{Afrika}
\author{Lengyel Zoltán Péter}
\date{Legutóbb frissítve: 2026. Május 12.}

\begin{document}

\maketitle
\tableofcontents
\newpage

\section{Afrika gazdasága}
\begin{itemize}
    \item Vízenergia
    \item Ásványkincsek (arany, ezüst, színesfémek, urán)
    \item Guineai-öböl, Észak-Afrika: kőolaj, földgáz
    \item Guineai-öböl partvidéke: kitermelt bauxit negyede
    \item Dél-Afrika: feketekőszén, gyémánt
    \item Marokkó: foszfát
    \item Kitermelt anyagok feldolgozás nélkül kerülnek a piacra (tőke és szakemberhiány miatt) 
    \item Nagy befektetők: Kína, India
    \item Turizmus: \begin{itemize}
        \item Egypitom - piramisok
        \item Zanzibár
        \item Kenya
    \end{itemize}
\end{itemize}

\section{Népesség Afrikában}
\begin{itemize}
    \item Népesség eloszlása egyeletlen (emberek nagyrésze a Nílus partján él)
    \item Nyomornegyedek
    \item Szahara és tőle északra: arab népesség iszlám vallással
    \item Szaharától délre: feketebőrűek
    \item Dél-Afrika: európai bevándorlók leszármazottjai is élnek
    \item Népesség 40\%-a 15 év alatti \begin{itemize}
        \item Magas a születések és elhalálozások száma 
        \item Alacsony az átlag életkor (kb. 20év) \begin{itemize}
            \item rossz egészségügy
            \item oktatás hiánya
            \item tiszta víz hiánya
            \item járványok
            \item magas csecsemőhalandóság
            \item háborúk
            \item éhínségek
        \end{itemize}
    \end{itemize}
\end{itemize}

\end{document}